\section{Methods}

\subsection{Pipeline Architecture}

GlycoSMILES2BAP operates through four sequential modules: Input Parsing, CCD Mapping, BAP Generation, and Output Formatting. The pipeline accepts IUPAC-condensed, WURCS, and GlycoCT input formats.

\subsection{CCD Mapper Module}

The CCD (Chemical Component Dictionary) provides standardized residue identifiers. Our mapper supports 28+ monosaccharide configurations. Key design decisions include: (1) case-insensitive matching, (2) anomeric position tracking (C2 for sialic acids, C1 for aldoses), and (3) ring oxygen positions (O4 for pentoses, O5 for hexoses, O6 for sialic acids).

\begin{table}[h]
\centering
\caption{Key CCD Mappings. The complete mapping table for 28+ monosaccharides is provided in Supplementary Table S1.}
\begin{tabular}{lllll}
\toprule
Monosaccharide & Anomer & Config & CCD Code & Anomeric C \\
\midrule
GlcNAc & $\beta$ & D & NAG & C1 \\
GlcNAc & $\alpha$ & D & A2G & C1 \\
Man & $\alpha$ & D & MAN & C1 \\
Man & $\beta$ & D & BMA & C1 \\
Gal & $\beta$ & D & GAL & C1 \\
Gal & $\alpha$ & D & GLA & C1 \\
Fuc & $\alpha$ & L & FUC & C1 \\
Neu5Ac & $\alpha$ & D & SIA & C2 \\
Neu5Gc & $\alpha$ & D & NGC & C2 \\
Xyl & $\beta$ & D & XYS & C1 \\
\bottomrule
\end{tabular}
\end{table}

\subsection{BAP Generator Module}

The bondedAtomPairs (BAP) generator converts glycan topology into AF3-compatible bond specifications. Each glycosidic linkage is represented as an explicit atom pair specifying donor residue, donor atom, acceptor residue, and acceptor atom.

\subsection{PDB Structure Validation Method}
\label{sec:pdb_validation}

To independently validate the CCD mapper's stereochemistry handling, we employed a PDB structure comparison approach. This validation was necessary because AlphaFold Server does not support direct glycan submission, and AF3 model parameters were not accessible.

\subsubsection{Validation Strategy}

We selected four PDB structures with known correct stereochemistry, each testing a distinct aspect of CCD mapping:

\begin{table}[h]
\centering
\caption{PDB Validation Test Cases}
\begin{tabular}{llll}
\toprule
PDB ID & Structure & Key Validation & Critical Test \\
\midrule
5NSC & Fucosylated & FUC = alpha-L-fucose & L-configuration \\
1L6R & Lactose & GAL = beta-D-galactose & Epimer distinction \\
2VXR & Sialyllactose & SIA anomeric = C2 & Ketose handling \\
5K65 & M3 N-glycan & Branch topology & Multi-residue \\
\bottomrule
\end{tabular}
\label{tab:pdb_cases}
\end{table}

\subsubsection{Critical Test: Sialic Acid C2 Anomeric Position}

The most critical validation case is PDB:2VXR (sialyllactose). Sialic acids (Neu5Ac) are ketoses with the anomeric carbon at C2, not C1. This is a known failure point for SMILES-based representations, which typically place glycosidic bonds at C1 by default. Successful CCD mapping must:
\begin{enumerate}
\item Map Neu5Ac(alpha) to SIA (not SLB for beta)
\item Assign anomeric position as C2 (not C1)
\item Assign ring oxygen as O6 (not O5)
\end{enumerate}

\subsection{Evaluation Metrics}
\label{sec:metrics}

We define three complementary accuracy metrics to comprehensively evaluate stereochemistry preservation. All metrics are calculated at the level of individual structural features (per monosaccharide residue or per glycosidic linkage).

\textbf{Epimer Accuracy} measures correct preservation of monosaccharide stereochemistry at all chiral centers. A residue is considered correctly identified if its CCD code matches the expected code, which encodes both monosaccharide type and absolute configuration (D/L).

\textbf{Anomeric Accuracy} measures correct anomeric configuration ($